\section{Where Is the Gap?}
\label{sec:gap}

Before designing a method, we measured on ImageNet-R (seed 1993) how much
accuracy is still reachable, and where the errors are.

\paragraph{Two reference points.}
Training jointly on all 200 classes, with the same recipe as task~0, gives
a final accuracy of $86.52$. The gap to RSIAT is therefore
$G=86.52-82.59=3.9$ points. If CA is given the \emph{oracle} statistics of
the old classes (means and covariances recomputed from their images with
the current network, which violates the exemplar-free constraint), AIA
rises by only $H=+0.195$ (two pairs, SD$(\Delta)=0.03$) and final accuracy
by $+0.5$. Feature drift is small: about $8\%$ of the feature norm per task.

\paragraph{The oldest classes, and confusion across tasks.}
Table~\ref{tab:gap} lists the final accuracy per task.

\begin{itemize}
  \item \emph{The oldest classes lose most.} Classes of tasks 0 and 1 lose
    $9.0$ and $8.3$ points against joint training, while the newest task is
    at the joint level.
  \item \emph{The excess errors are cross-task.} Per pair of classes, the
    ratio of cross-task to within-task confusions is $1.38$ for RSIAT and
    $0.90$ for joint training.
  \item \emph{Oracle statistics move the errors instead of removing them.}
    They lift task~0 from $80.1$ to $87.1$, but tasks 5--9 drop by $1$--$2.6$
    points, because newer classes are now predicted as task~0 more often.
\end{itemize}

The remaining gap therefore lies in how the classes of different tasks
overlap in the learned representation, and not in the accuracy of the
statistics. Inference-time changes can only redistribute this overlap
better. The next section does this with information the linear head
ignores.

\begin{table}[t]
\centering
\caption{Final accuracy (\%) per task on ImageNet-R (task~0 holds the oldest
classes). CIL: RSIAT, mean of 4 replicates; oracle: CA with oracle statistics
of old classes, 2 runs; joint: one task with all 200 classes.}
\label{tab:gap}
\setlength{\tabcolsep}{4pt}
\begin{adjustbox}{max width=\textwidth}
\begin{tabular}{lrrrrrrrrrr}
\toprule
Task & 0 & 1 & 2 & 3 & 4 & 5 & 6 & 7 & 8 & 9\\
\midrule
Joint  & 89.1 & 88.7 & 86.4 & 86.4 & 83.2 & 86.0 & 85.8 & 87.0 & 85.4 & 85.8\\
CIL    & 80.1 & 80.4 & 82.7 & 82.3 & 78.8 & 82.7 & 85.3 & 84.2 & 83.9 & 86.2\\
Oracle & 87.1 & 83.1 & 83.5 & 82.2 & 79.1 & 81.6 & 84.7 & 82.9 & 82.1 & 83.6\\
\midrule
Joint $-$ CIL & $+9.0$ & $+8.3$ & $+3.7$ & $+4.1$ & $+4.4$ & $+3.3$ & $+0.5$ & $+2.9$ & $+1.6$ & $-0.3$\\
\bottomrule
\end{tabular}
\end{adjustbox}
\end{table}

